\subsubsection{Jump Marker}
\label{SEC:JumpMarker}
\key*{}[marker!jump marker]

%%%%%%%%%%%%%%%%%%%%%%%%%%%%%%%%%%%%%%%%%%%%%%%%%%%%%%%%%%%%%%%%%%%%%%%%%%%%%%%%

%% Note:
%% - Krivodonova_2004: Original jump marker paper (as far as I can tell)
%% - Burgess_2011aiaa: Provides the jump marker with    log10 scaling.
%% - Burgess_2011thesis: Burgess's PhD that is similar to Burgess_2011aiaa.
%% - Barter_2007:  Provides the jump marker without log10 scaling
%%                 and coefficients (C1,C2,dPsi) for the scaled marker.
%% - Barter_2008:  Barter's PhD that is similar to Barter_2007.

For flows with shocks or other sudden discontinuities,
it is useful to mark those cells that experience a sudden jump
in pressure or density.
%
The jump marker\cite{Krivodonova_2004, Barter_2007, Barter_2008,
Burgess_2011aiaa, Burgess_2011thesis} is a method for identifying such sudden
changes across adjacent cells.
%
Using the notation in \fref{FIG:Jump_Marker}, the jump indicator
can be defined as:
\begin{align}
  \label{EQN:JumpIndicator1}
    J_\kappa
 =& \frac{1}{\|\partial \kappa\|}
    \int_{\partial \kappa}
       \left|
          \frac{ [[ \psi ]] \cdot \vec{n} }
               { \{ \psi \}               }
       \right|
       ds
  & \begin{aligned}
      [[ \psi ]] =& \psi_1 \vec{n}_1 + \psi_2 \vec{n}_2
                 =  \left( \psi_1 - \psi_2 \right) \vec{n}_1 \\
      \{ \psi \} =& \left( \psi_1 + \psi_2 \right) / 2
    \end{aligned}
 \\
  \label{EQN:JumpIndicator2}
 =& \frac{1}{\|\partial \kappa\|}
    \int_{\partial \kappa}
       \left|
         \frac{ \left( \psi_1 - \psi_2 \right) }
              { \left( \psi_1 + \psi_2 \right) / 2 }
       \right|
       ds
 =  \left[ 0 : 2 \right]
\end{align}
%
Here,
$[[\psi]]$ represents the jump across adjacent face values and
$\{\psi\}$ represents the average of those values.
%
The jump indicator $J_\kappa$ can vary between 0 (no jump)
and 2 (infinite jump).

\begin{figure}
  \begin{centering}
    \insertimage[width=0.5\textwidth]{tikz/Jump_Marker}
   %\input{Figures/Jump_Marker}
    \caption{Jump Marker Notation}
    \label{FIG:Jump_Marker}
  \end{centering}
\end{figure}

The jump marker is obtained by taking the logarithm of the value.
\begin{align}
    {\cal M} \left( J_\kappa \right)
 =& \log_{10} \left[ \max \left( J_\kappa , 2^{-23} \right) \right]
 =  \left[ -6.9 : 0.3 \right]
  \label{EQN:JumpMarker}
\end{align}
The jump marker ${\cal M}$ can vary between -6.9 (no jump)
and 0.3 (infinite jump).
%
The scaled marker is then computed as described in \sref{SEC:ScaledMarker}.

%%%%%%%%%%%%%%%%%%%%%%%%%%%%%%%%%%%%%%%%%%%%%%%%%%%%%%%%%%%%%%%%%%%%%%%%%%%%%%%%

\SeeAlso{
\inp[nmlLimiter::MarkerOption]{Marker\_Option},
\inp[nmlLimiter::MarkerVariable]{Marker\_Variable}
}

%%%%%%%%%%%%%%%%%%%%%%%%%%%%%%%%%%%%%%%%%%%%%%%%%%%%%%%%%%%%%%%%%%%%%%%%%%%%%%%%

